% =====================================================================
% part2_mpi.tex -- PARTE 2: Estensione multi-simulazione con MPI
% Membro responsabile: ___________________
% Puoi modificare liberamente da qui in poi, non toccare gli altri file.
% =====================================================================
\section{Part 2 -- Multi-Simulation Extension with MPI}
\label{sec:mpi}

\subsection{Problem Description}
The second part of the project extends the OpenMP simulator of Section~\ref{sec:openmp} to
the distributed-memory setting using MPI, with the goal of executing \emph{three independent
simulations simultaneously} on an HPC cluster. Each simulation continues to exploit
intra-node, shared-memory parallelism through OpenMP, resulting in a hybrid MPI+OpenMP
implementation. The three simulations differ in their initial conditions:

\begin{enumerate}[label=(\roman*)]
    \item \textbf{Simulation 1}: identical to the base case of Part~1, with a single impulse
    source at $(i_0,j_0)$ and $t=0$.
    \item \textbf{Simulation 2}: two independent wave sources, located at two distinct points
    of the plane, sharing the same damping $\gamma$ and propagation speed $c$, both activated
    at $t=0$. This configuration is designed to highlight wave interference patterns arising
    from the superposition of the two wavefronts.
    \item \textbf{Simulation 3}: identical to Simulation 2, but the second source is activated
    at a later time $t = t_{\text{start}} > 0$, producing an asymmetric interference pattern.
\end{enumerate}

\subsection{Parallelization Strategy}
The three simulations are mapped onto separate MPI processes (or groups of processes),
allowing them to run concurrently and independently, since there is no data dependency between
them. Each MPI process executes its own instance of the OpenMP-parallelized time-stepping loop
described in Section~\ref{sec:openmp}, writing its output frames to a dedicated folder
(\texttt{sim1}, \texttt{sim2}, \texttt{sim3} respectively). This two-level parallelization
scheme -- MPI across simulations and OpenMP within each simulation -- allows efficient use of
a multi-node cluster where each node hosts one MPI process exploiting all of its available
cores via OpenMP threads.

% TODO (optional): describe communication pattern (if any communication between
% processes is needed, e.g., for synchronized output or load balancing), and how
% the two interfering sources are summed/composed onto the shared grid in
% Simulations 2 and 3.

\subsection{Results}
\textit{[Insert here: simulation parameters for the three runs (Table~\ref{tab:part2-params}),
example frames illustrating single-source propagation and the interference patterns of
Simulations 2 and 3 (Figure~\ref{fig:part2-frames}), and the measured execution time/scalability
of the hybrid MPI+OpenMP implementation as a function of the number of MPI processes and OpenMP
threads.]}

\begin{table}[h]
\centering
\caption{Simulation parameters -- Part 2.}
\label{tab:part2-params}
\begin{tabular}{lccc}
\toprule
Parameter & Sim. 1 & Sim. 2 & Sim. 3 \\
\midrule
$\gamma$               & \textit{TBD} & \textit{TBD} & \textit{TBD} \\
$c$                    & \textit{TBD} & \textit{TBD} & \textit{TBD} \\
Source(s) position     & \textit{TBD} & \textit{TBD} & \textit{TBD} \\
$t_{\text{start}}$ (2nd source) & -- & 0 & \textit{TBD} \\
\bottomrule
\end{tabular}
\end{table}

\begin{figure}[h]
\centering
% \includegraphics[width=\linewidth]{frames_mpi.png}
\caption{Comparison of single-source propagation and multi-source interference patterns. \textit{[Placeholder -- insert figure]}}
\label{fig:part2-frames}
\end{figure}

\begin{figure}[h]
\centering
% \includegraphics[width=\linewidth]{scaling_mpi.png}
\caption{Execution time and scalability of the hybrid MPI+OpenMP implementation. \textit{[Placeholder -- insert figure]}}
\label{fig:part2-scaling}
\end{figure}
